
\label{sec:ablation_vfm}

To demonstrate the effectiveness of our chosen combination of DINOv2~\cite{oquab2023dinov2} and DPAv3 as VFM targets, we evaluate a diverse set of alternative VFMs on the CALVIN~\cite{mees2022calvin} benchmark. We categorize them into semantic, geometric, motion-aware, and synergistic targets. The results are detailed in \cref{tab:ablation_vfm}.

\noindent \textbf{Semantic Targets (Dense vs. Global).} By extracting dense patch-level representations,
DINOv2~\cite{oquab2023dinov2} and DINOv3~\cite{simeoni2025dinov3} significantly outperform global semantic models CLIP~\cite{radford2021learning} and SigLIP~\cite{zhai2023sigmoid} (e.g., 4.01 for DINOv2 vs. 3.72 for CLIP in average length). This gap demonstrates that while global semantics provide scene-level context, precise manipulation inherently relies on dense semantic distinctiveness for the action expert to attend to task-relevant affordances. 

\noindent \textbf{Geometric Targets (Static vs. Dynamic).} 
DPAv3~\cite{lin2025depth} significantly outperforms VGGT~\cite{wang2025vggt} (3.99 vs. 3.74). We attribute this performance margin to DPAv3 being trained to adapt to dynamic video streams, enabling it to extract dynamic geometric structures. Conversely, VGGT focuses on static scene reconstruction, which makes it fundamentally unsuited to handle the dynamic object interactions inherent in robotic tasks.

\noindent \textbf{Motion-aware Targets.} 
We also explore using the representations extracted by video foundation models, VideoMAEv2~\cite{wang2023videomae} and V-JEPA2~\cite{assran2025v}, as motion-aware targets. However, they demonstrate sub-optimal results. We hypothesize that current video models still lag behind specialized image counterparts (e.g., DINOv2) in fine-grained feature fidelity.

\noindent \textbf{Synergistic Targets.} 
The best performance is achieved by fusing complementary modalities. Our choice (DINOv2 \& DPAv3) achieves the highest average sequence length of 4.16. Replacing DPAv3 with VGGT (DINOv2 \& VGGT) results in a performance drop to 4.04, confirming that static-scene geometric priors are insufficient for dynamic manipulation tasks.
Second, substituting DINOv2 with SigLIP (SigLIP \& DPAv3) leads to a decline to 4.06, demonstrating that dense and patch-level representations are more critical for low-level control than global-level language-aligned features.

\begin{table}[t]
\centering
\small
\caption{Quantitative comparison of different VFM representations as teacher targets.}
\label{tab:ablation_vfm}

\definecolor{mygray}{gray}{0.95}

\renewcommand{\arraystretch}{1.00}
\scalebox{0.96}{
\begin{tabular}{c|c|ccccc|c}
\toprule
\multirow{2}{*}{\textbf{Type}} & \multirow{2}{*}{\shortstack{\textbf{Representation} \\ \textbf{Source}}} & \multicolumn{5}{c|}{\textbf{$i^{th}$ Task Success Rate}} & \multirow{2}{*}{\textbf{Avg. Len.} $\uparrow$} \\
\cmidrule(lr){3-7}
 & & \textbf{1} & \textbf{2} & \textbf{3} & \textbf{4} & \textbf{5} & \\
\midrule


\multirow{4}{*}{\shortstack{Semantic \\ Targets}} 
 & CLIP~\cite{radford2021learning}      & 94.1 & 82.0 & 72.5 & 65.0 & 58.0 & 3.72 \\
 & SigLIP~\cite{zhai2023sigmoid}               & 94.5 & 84.8 & 74.7 & 65.5 & 57.0 & 3.77 \\
 & DINOv2~\cite{oquab2023dinov2}  & 94.1 & 87.1 & 79.3 & 73.5 & 66.5 & 4.01 \\
 & DINOv3~\cite{simeoni2025dinov3}               & 95.6 & 87.7 & 78.8 & 70.3 & 62.1 & 3.95 \\
\midrule


\multirow{2}{*}{\shortstack{Geometric \\ Targets}} 
 & DPAv3~\cite{lin2025depth} & 94.0 & 87.1 & 80.4 & 73.2 & 64.1 & 3.99 \\
 & VGGT~\cite{wang2025vggt}   & 93.5 & 84.3 & 74.8 & 65.6 & 55.8 & 3.74 \\ 

\midrule

\multirow{2}{*}{\shortstack{Motion-aware \\ Targets}} 
 & VideoMAEv2~\cite{wang2023videomae}  & 94.6 & 85.9 & 78.0 & 70.2 & 60.8 & 3.90 \\
 & V-JEPA2~\cite{assran2025v}  & 93.1 & 84.3 & 74.3 & 65.6 & 56.4 & 3.74 \\
\midrule


  & SigLIP \& DPAv3       & \textbf{96.2} & 88.3 & 80.7 & 74.2 & 66.1 & 4.06 \\
  & DINOv2 \& VGGT     & 95.7 & 88.8 & 80.9 & 73.9 & 64.9 & 4.04 \\

\rowcolor{mygray}
\cellcolor{white} \multirow{-3}{*}{\shortstack{Synergistic \\ Targets}} & \textbf{DINOv2 \& DPAv3 (our)} & {95.8} & {\textbf{90.7}} & {\textbf{83.7}} & {\textbf{77.0}} & {\textbf{68.9}} & {\textbf{4.16}} \\
\bottomrule
\end{tabular}
}
\end{table}

